\section{Kapillarlik hodisasi, kapillar nay}



\begin{enumerate}

\item Agar suyuqlik molekulalari orasidagi o‘zaro ta’sir kuchi suyuqlik molekulasi bilan qattiq jism molekulasi orasidagi o‘zaro ta’sir kuchidan kichik bo‘lsa, qanday hodisa ro‘y beradi?
\begin{enumerate}[label=\Alph*)]
\item Dinamik muvozanat
\item Qaynash
\item Sublimatsiya
\item Xo‘llash
\end{enumerate}

\item Agar suyuqlik molekulalari orasidagi o‘zaro ta’sir kuchi suyuqlik molekulasi bilan qattiq jism molekulasi orasidagi o‘zaro ta’sir kuchidan katta bo‘lsa, qanday hodisa ro‘y beradi?
\begin{enumerate}[label=\Alph*)]
\item Dinamik muvozanat
\item Qaynash
\item Xo‘llamaslik
\item Xo‘llash
\end{enumerate}

\item Xo‘llaydigan suyuqlik uchun xo‘llash burchagi necha gradus bo‘lishi mumkin?
\begin{enumerate}[label=\Alph*)]
\item 67
\item Barchasi to‘g‘ri
\item 41
\item 12

\end{enumerate}

\item Xo‘llaydigan suyuqlik uchun xo‘llash burchagi quyidagilardan qaysi biri bo‘lishi mumkin?
\begin{enumerate}[label=\Alph*)]
\item 67
\item 112
\item 90
\item A va B javoblar to‘g‘ri
\end{enumerate}

\item Xo‘llamaydigan suyuqlik uchun xo‘llash burchagi necha gradus bo‘lishi mumkin?
\begin{enumerate}[label=\Alph*)]
\item 105
\item 150
\item 71
\item A va B javoblar to‘g‘ri
\end{enumerate}


\item Radiusi 24 $\mu m$ bo‘lgan kapillyar naychada kerosin qanday balandlikka ko‘tariladi (sm)? Kerosinning sirt taranglik koeffitsiyenti 24 mN/m ga, $\rho = 0,8 \, g/sm^3$ ga teng.
\begin{enumerate}[label=\Alph*)]
\item 25
\item 12,5
\item 50
\item 75
\end{enumerate}

\item Diametri 0,3 mm bo‘lgan kapillyarda kerosin 20 mm ko‘tarildi. Kerosinning sirt taranglik koeffitsiyentini aniqlang (mN/m). $\rho = 0,8 \, g/sm^3$.
\begin{enumerate}[label=\Alph*)]
\item 8
\item 10
\item 12
\item 14
\end{enumerate}

\item Diametri 0,73 mm bo‘lgan kapillyar nayda suv qancha balandlikka (sm) ko‘tariladi? Suvning sirt taranglik koeffitsiyenti 73 mN/m ga teng.
\begin{enumerate}[label=\Alph*)]
\item 2
\item 1
\item 4
\item 6
\end{enumerate}

\item Yerda ho‘llovchi suyuqlik kapillyar bo‘ylab 4 sm ko‘tariladi. Erkin tushish tezlanishi 5 $m/s^2$ bo‘lgan sayyoralarda bu balandlik qanday bo‘ladi (sm)? Yerda $g = 10 \, m/s^2$.
\begin{enumerate}[label=\Alph*)]
\item 8
\item 2
\item 4
\item 1
\end{enumerate}

\item Suyuqlikning kapillyar naychadagi ko‘tarilish balandligi 10 mm. Radiusi 3 marta kichik bo‘lgan boshqa kapillyar naychada shu suyuqlik qanday balandlikka ko‘tariladi?
\begin{enumerate}[label=\Alph*)]
\item 3 m
\item 3 sm
\item 0,3 m
\item 30 mm
\end{enumerate}

\item Yer sirtida joylashgan kapillyar naychada suv 24 mm ga ko‘tariladi. Oyda erkin tushish tezlanishi Yerdagidan 6 marta kichik. Oyda shu naychadagi suv qanday (mm) balandlikka ko‘tariladi?
\begin{enumerate}[label=\Alph*)]
\item 134
\item 35
\item 144
\item 102
\end{enumerate}

\item Ichki diametrlari 1 va 2 mm bo‘lgan ikkita kapillyardagi suv sathlarining farqini aniqlang (mm). $\sigma = 72 \, mN/m$.
\begin{enumerate}[label=\Alph*)]
\item 0
\item 14,4
\item 28,8
\item 43,2
\end{enumerate}

\item Ikkita kapillyar nay ikkita suyuqlikka tushirilgan. 2-nayning radiusi 4 marta kichik va u sirt taranglik koeffitsiyenti 2 marta katta bo‘lgan suyuqlikka botirilgan. 2-nayda suyuqlik 1-naydagidan necha marta baland ko‘tariladi?
\begin{enumerate}[label=\Alph*)]
\item 2
\item 4
\item 8
\item 16
\end{enumerate}

\item Ikkala uchi ochiq bo‘lgan shisha kapillyar pastki uchi bilan suvga tushirildi. ($\sigma = 73 \, mN/m$). Agar suv shishani to‘liq ho‘llasa suv kapillyarda qanday balandlikka (mm) ko‘tariladi? Kapillyarning ichki radiusi 1 mm ga teng.
\begin{enumerate}[label=\Alph*)]
\item 7,3
\item 14,6
\item 21,9
\item 29,2
\end{enumerate}

\item Simobli barometr naychasining diametri 3 mm. Agar simobning kapillyar pasayishini hisobga olsak, barometrning ko‘rsatishiga qanday tuzatish kiritish lozim (mm)? $\sigma = 22 \, mN/m$.
\begin{enumerate}[label=\Alph*)]
\item 5,9
\item 5,1
\item 4,7
\item 4,2
\end{enumerate}

\item Uchta bir xil kapillyar nay bir uchi bilan ularni ho‘llovchi va sirt taranglik koeffitsiyentlari $\sigma, 2\sigma$ va $3\sigma$ bo‘lgan suyuqliklarga tushirilgan. Suyuqlikning ko‘tarilish balandliklari orasidagi munosabatni aniqlang.
\begin{enumerate}[label=\Alph*)]
\item $h_1 = h_2 = h_3$
\item $h_1 = h_2 = 3h_3$
\item aniqlab bo‘lmaydi, chunki zichliklar berilmagan
\item $h_1 = h_2/2 = h_3/3$
\end{enumerate}

\item U-simon kapillyar nayning o‘ng va chap tirsaklarining ichki diametrlari mos ravishda 1 mm va 0,2 mm ga teng. Nayga quyilgan suv sathlari farqini toping (sm). $\sigma = 73 \, mN/m$.
\begin{enumerate}[label=\Alph*)]
\item 1,2
\item 1,46
\item 18
\item 2,4
\end{enumerate}

\item Ikkita parallel plastina va kapillyar nay ho‘llaydigan suyuqlikka tushirildi. Plastinalar orasidagi masofa kapillyar nayning radiusiga teng bo‘lsa, ularda ko‘tarilgan suyuqlik ustunlarining balandliklari $h_1$ va $h_2$ qanday munosabatda bo‘ladi?
\begin{enumerate}[label=\Alph*)]
\item $h_1 = 1,5h_2$
\item $2h_1 = h_2$
\item $h_1 = 2h_2$
\item $h_1 = h_2$
\end{enumerate}

\item Oralaridagi masofasi 0,73 mm bo‘lgan parallel plastinkada suv qancha balandlikka ko‘tariladi (sm)? Suvning sirt taranglik koeffitsiyenti 73 mN/m ga teng.
\begin{enumerate}[label=\Alph*)]
\item 2
\item 1
\item 4
\item 6
\end{enumerate}

\item Oralaridagi masofasi 2 mm bo‘lgan parallel plastinkada suv qancha balandlikka (sm) ko‘tariladi? Suvning sirt taranglik koeffitsiyenti 73 mN/m ga teng.
\begin{enumerate}[label=\Alph*)]
\item 7,3
\item 1,76
\item 0,73
\item 0,365
\end{enumerate}

\item Parallel plastinkalar oralaridagi masofa necha mm bo‘lganda, suv 7,3 mm ga ko‘tariladi? Suvning sirt taranglik koeffitsiyenti 73 mN/m ga teng.
\begin{enumerate}[label=\Alph*)]
\item 7,3
\item 1,76
\item 2
\item 0,365
\end{enumerate}

\item Kapillyar nayning radiusi 0,5 mm va undagi suyuqlikning sirt taranglik koeffitsiyenti 40 mN/m bo‘lsa, qo‘-shimcha bosimni aniqlang (Pa).
\begin{enumerate}[label=\Alph*)]
\item 80
\item 160
\item 240
\item 320
\end{enumerate}

\end{enumerate}




\begin{table}[h!]
\centering
\caption{Xo‘llash va Kapillyar hodisalar bo‘limlari uchun to‘g‘ri javoblar jadvali (1--22)}
\begin{tabular}{|c|c|c|c|c|c|}
\hline
\textbf{№} & \textbf{Javob} & \textbf{№} & \textbf{Javob} & \textbf{№} & \textbf{Javob} \\ \hline
1 & D & 9 & A & 17 & B \\ \hline
2 & C & 10 & D & 18 & D \\ \hline
3 & B & 11 & C & 19 & A \\ \hline
4 & D & 12 & B & 20 & C \\ \hline
5 & D & 13 & C & 21 & C \\ \hline
6 & A & 14 & B & 22 & B \\ \hline
7 & C & 15 & B & & \\ \hline
8 & C & 16 & C & & \\ \hline
\end{tabular}
\end{table}
